% SPDX-License-Identifier: GPL-3.0-or-later
% Copyright (C) 2026 Alessandro Greco
% Edición en español internacional. Conserva la infraestructura técnica compartida.
\usepackage{fontspec}
\usepackage{polyglossia}
\setmainlanguage{spanish}
\setotherlanguage{italian}
\setotherlanguage{english}
\usepackage[autostyle=once]{csquotes}
\setmainfont{TeX Gyre Pagella}
\setsansfont{TeX Gyre Heros}[Scale=MatchLowercase]
\setmonofont{DejaVu Sans Mono}[Scale=0.85]
\usepackage[a4paper,inner=3cm,outer=2.5cm,top=2.5cm,bottom=3cm]{geometry}
\usepackage{microtype}
\setlength{\emergencystretch}{3em}
\usepackage{xcolor}
\usepackage{graphicx}
\usepackage{booktabs,tabularx,longtable,multirow}
\usepackage{amsmath,amssymb}
\usepackage{siunitx}
% Coma decimal y espacio fino, sin depender de un perfil locale=ES inexistente.
\sisetup{output-decimal-marker={,},group-separator={\,}}
\usepackage{enumitem}
\usepackage{tikz}
\usetikzlibrary{arrows.meta,positioning,shapes,calc,fit,backgrounds,decorations.pathreplacing,decorations.pathmorphing,matrix,chains}
\usepackage{bytefield}
\usepackage{listings}
\usepackage[most]{tcolorbox}
\usepackage[backend=biber,style=alphabetic,maxbibnames=99]{biblatex}
\addbibresource{../shared/bibliography/tor.bib}
% La bibliografía compartida está escrita para la edición italiana: aquí se sustituyen las notas
% en italiano. «Consultato il ...» pasa a url + urldate, que biblatex escribe en este idioma.
\DeclareSourcemap{\maps[datatype=bibtex,overwrite=true]{
  \map{
    \step[fieldsource=note, match={Consultato il 23 settembre 2026}, final]
    \step[fieldsource=howpublished, fieldtarget=url]
    \step[fieldsource=url, match=\regexp{\A\\url\{(.+)\}\z}, replace={$1}]
    \step[fieldset=urldate, fieldvalue={2026-09-23}]
    \step[fieldset=note, null]
  }
  \map{\step[fieldsource=entrykey, match=\regexp{\Atorblogwebtunnel\z}, final]\step[fieldset=note, fieldvalue={12 de marzo de 2024}]}
  \map{\step[fieldsource=entrykey, match=\regexp{\Aperry2018torbrowser\z}, final]\step[fieldset=note, fieldvalue={Borrador}]}
  \map{\step[fieldsource=entrykey, match=\regexp{\Adsa2022\z}, final]\step[fieldset=note, fieldvalue={DO L 277 de 27.10.2022}]}
  \map{\step[fieldsource=entrykey, match=\regexp{\Acodiceprivacy\z}, final]\step[fieldset=note, fieldvalue={Texto consolidado publicado por la autoridad italiana de protección de datos (Garante per la protezione dei dati personali)}]}
  \map{\step[fieldsource=entrykey, match=\regexp{\Awinter2014spoiled\z}, final]\step[fieldset=note, fieldvalue={incorporado después, con otros autores, en \emph{Privacy Enhancing Technologies Symposium} (PETS 2014)}]}
  \map{\step[fieldsource=entrykey, match=\regexp{\Alegge167\z}, final]\step[fieldset=note, fieldvalue={Gazzetta Ufficiale, Serie Generale n.º~277, 27.11.2017}]}
  \map{\step[fieldsource=entrykey, match=\regexp{\Anistpqc2024\z}, final]\step[fieldset=note, fieldvalue={los tres estándares poscuánticos, publicados conjuntamente}]}
  \map{\step[fieldsource=entrykey, match=\regexp{\Ahayden2014metadata\z}, final]\step[fieldset=note, fieldvalue={en el Johns Hopkins Foreign Affairs Symposium, 1 de abril de 2014}]}
  \map{\step[fieldsource=entrykey, match=\regexp{\Apfitzmann2010terminology\z}, final]\step[fieldset=howpublished, fieldvalue={Versión 0.34}]}
}}
\DeclareBibliographyAlias{legislation}{misc}
\usepackage{imakeidx}
\makeindex[intoc,title={Índice analítico}]
\usepackage[hidelinks]{hyperref}
\usepackage[nameinlink]{cleveref}
\crefname{appendice}{apéndice}{apéndices}
\Crefname{appendice}{Apéndice}{Apéndices}
